\documentclass[10pt,twocolumn]{article}

\usepackage[margin=1.9cm,top=2.2cm,bottom=2.2cm]{geometry}
\usepackage{amsmath}
\usepackage{amssymb}
\usepackage{hyperref}
\usepackage{times}
\usepackage{titlesec}
\usepackage{booktabs}
\usepackage{microtype}
\usepackage{parskip}
\hypersetup{colorlinks=true, linkcolor=blue, citecolor=blue, urlcolor=blue}

\titleformat{\section}{\normalfont\bfseries\small\centering}{}{0em}{}
\titlespacing{\section}{0pt}{6pt}{3pt}

\setlength{\columnsep}{0.5cm}
\setlength{\parskip}{3pt}
\setlength{\parindent}{1em}

\usepackage{fancyhdr}
\pagestyle{fancy}
\fancyhf{}
\fancyhead[L]{\small\itshape BCT Superfluid Lattice Model --- Letter 25}
\fancyhead[R]{\small\itshape March 2026}
\fancyfoot[C]{\thepage}
\renewcommand{\headrulewidth}{0.4pt}

\begin{document}

\twocolumn[{%
\begin{center}
{\large\bfseries The Weinberg Angle to Sub-0.01\% from BCT Oblique Corrections:}\\[3pt]
{\large\bfseries $\sin^2\theta_W = r_{\rm tet}^2/(r_{\rm oct}^2 + r_{\rm tet}^2)\,(1 + 9\alpha_0/4)$}\\[8pt]
{\normalsize Michel Robert Cabri\'{e}$^{1,\,*}$}\\[2pt]
{\small $^1$\textit{Independent Researcher;  ORCID: \href{https://orcid.org/0009-0007-9561-9859}{0009-0007-9561-9859}}}\\[2pt]
{\small (Dated: March 2026)\quad BCT Letter 25}\\[6pt]
\end{center}
\begin{quotation}
\small\noindent
The weak mixing angle $\sin^2\theta_W$ is derived from BCT void geometry. The tree-level formula
$\sin^2\theta_W^{(0)} = r_{\rm tet}^2/(r_{\rm oct}^2 + r_{\rm tet}^2) = 0.22744$ gives $-1.6\%$
below the PDG on-shell value. The NLO oblique correction with universal BCT coefficient $N = 9/4$
yields:
\[
  \sin^2\theta_W = \sin^2\theta_W^{(0)}\!\left(1 + \frac{9\alpha_0}{4}\right) = 0.231229\,,
\]
in agreement with the PDG value $0.23122(3)$ at $+0.004\%$ (Prediction \#99). The tree-level
formula reveals that the Weinberg angle is the ratio of tetrahedral to total void areas in the
BCT unit cell.
\end{quotation}
\smallskip
{\small PACS: 12.15.Lk, 11.30.Qc, 12.10.-g, 12.15.-y}
\vspace{6pt}
}]

\section{Introduction}

In the BCT Superfluid Lattice Model~[3,\,4], electroweak symmetry breaking arises from the BCT
void topology: the octahedral void hosts $\mathrm{SU}(2)_L$ gauge fields and the tetrahedral void
hosts $\mathrm{U}(1)_Y$. The Weinberg angle is the ratio of these void coupling strengths.

\section{Tree-Level Weinberg Angle}

The geometric tree-level formula is:
\begin{align}
\sin^2\theta_W^{(0)}
  &= \frac{r_{\rm tet}^2}{r_{\rm oct}^2 + r_{\rm tet}^2} \notag\\
  &= \frac{(0.11237)^2}{(0.20711)^2 + (0.11237)^2}
   = 0.22744\,.
\label{eq:tree}
\end{align}
This is $-1.6\%$ below the PDG on-shell value, requiring an NLO correction.

\section{NLO Oblique Correction}

BCT NLO oblique corrections follow $O^{\rm NLO} = O^{(0)}(1 + N\alpha_0)$. For the Weinberg
angle, the universal coefficient is $N = 9/4$:
\begin{align}
\sin^2\theta_W
  &= \sin^2\theta_W^{(0)}\!\left(1 + \frac{9\alpha_0}{4}\right) \notag\\
  &= 0.22744 \times 1.016668
   = 0.231229\,.
\label{eq:NLO}
\end{align}

\section{Prediction \#99}

\begin{table}[h]
\centering
\small
\begin{tabular}{@{}lcc@{}}
\toprule
 & BCT & PDG~[1] \\
\midrule
$\sin^2\theta_W^{(0)}$ (tree) & 0.22744 & 0.23122(3) \\
$\sin^2\theta_W^{(\rm NLO)}$  & 0.231229 & \\
Error & $+0.004\%$ & \\
\bottomrule
\end{tabular}
\end{table}

The BCT programme now counts 99 sub-1\% predictions from zero free parameters. The residual
$+0.004\%$ ($+0.000009$ in $\sin^2\theta_W$) is at the level of two-loop electroweak corrections,
within current experimental precision $\pm0.00003$.

\section{The 9/4 Universality Class}

The coefficient $N = 9/4$ belongs to a small set of universal BCT oblique coefficients:
\begin{itemize}\setlength\itemsep{1pt}
\item $N = 3/8$ --- deconfinement temperature $T_c$ (Letter~24 [6]); top quark NLO.
\item $N = 9/4$ --- Weinberg angle; $\tau$ anomalous magnetic moment.
\item $N = 4\alpha_0/\pi$ --- lepton hierarchy self-consistency (App.~AB).
\end{itemize}
The ratio $(9/4)/(3/8) = 6 = N_c!$ may reflect the structure of the BCT loop measure.

\section{Conclusion}

\begin{align}
\sin^2\theta_W
  &= \frac{r_{\rm tet}^2}{r_{\rm oct}^2 + r_{\rm tet}^2}
    \left(1 + \frac{9\alpha_0}{4}\right) \notag\\
  &= 0.231229\quad(+0.004\%).
\label{eq:final}
\end{align}

Prediction \#99 from zero free parameters. The Weinberg angle is a direct consequence of the
BCT void geometry.

\bigskip
\noindent$^*$ \href{mailto:ZeroFreeParameters@gmail.com}{\texttt{ZeroFreeParameters@gmail.com}}

\begin{thebibliography}{6}

\bibitem{PDG2022}
R.~L. Workman et al.\ (PDG), Prog.\ Theor.\ Exp.\ Phys.\ \textbf{2022}, 083C01 (2022).

\bibitem{ErlerFreitas2024}
J.~Erler and A.~Freitas, PDG Electroweak Model Chapter (2024).

\bibitem{BCT_PRL}
M.~R. Cabri\'{e}, ``Zero Free Parameters\ldots'' \href{https://doi.org/10.5281/zenodo.18884415}{DOI:~10.5281/zenodo.18884415} (2026).

\bibitem{BCT_Monograph}
M.~R. Cabri\'{e}, ``The BCT Superfluid Lattice Model,'' \href{https://doi.org/10.5281/zenodo.18884976}{DOI:~10.5281/zenodo.18884976} (2026).

\bibitem{BCT_StrongCP}
M.~R. Cabri\'{e}, ``Geometric Resolution of the Strong CP Problem without an Axion,''
\href{https://doi.org/10.5281/zenodo.18885628}{DOI:~10.5281/zenodo.18885628} (2026).

\bibitem{BCT_L24}
M.~R. Cabri\'{e}, BCT Letter~24 (2026).

\end{thebibliography}

\end{document}
